% !TEX TS-program = XeLaTeX
% use the following command:
% all document files must be coded in UTF-8
\documentclass[portuguese]{textolivre}
% build HTML with: make4ht -e build.lua -c textolivre.cfg -x -u article "fn-in,svg,pic-align"

\journalname{Texto Livre}
\thevolume{19}
%\thenumber{1} % old template
\theyear{2026}
\receiveddate{\DTMdisplaydate{2026}{2}{12}{-1}} % YYYY MM DD
\accepteddate{\DTMdisplaydate{2026}{2}{27}{-1}}
\publisheddate{\DTMdisplaydate{2026}{3}{18}{-1}}
\corrauthor{Lucas Jansen}
\articledoi{10.1590/1983-3652.2026.64400}
%\articleid{NNNN} % if the article ID is not the last 5 numbers of its DOI, provide it using \articleid{} commmand 
% list of available sesscions in the journal: articles, dossier, reports, essays, reviews, interviews, editorial
\articlesessionname{reviews}
\runningauthor{Jansen} 
%\editorname{Leonardo Araújo} % old template
\sectioneditorname{Daniervelin Pereira~\orcid{0000-0003-1861-3609}}
\layouteditorname{Leonardo Araújo~\orcid{0000-0003-3884-2177}}

\title{Resenha de A geração ansiosa: como a infância hiperconectada está causando uma epidemia de transtornos mentais}
\othertitle{Review of The anxious generation: how the great rewiring of childhood is causing an epidemic of mental illness}
% if there is a third language title, add here:
%\othertitle{Artikelvorlage zur Einreichung beim Texto Livre Journal}

\author[1]{Lucas Jansen~\orcid{0000-0003-4110-8117}\thanks{Email: \href{mailto: lucaslimajansen94@gmail.com}{lucaslimajansen94@gmail.com}}}
\affil[1]{Universidade de Brasília, Brasília, DF, Brasil.}

\addbibresource{article.bib}
% use biber instead of bibtex
% $ biber article

% used to create dummy text for the template file
\definecolor{dark-gray}{gray}{0.35} % color used to display dummy texts
\usepackage{lipsum}
\SetLipsumParListSurrounders{\colorlet{oldcolor}{.}\color{dark-gray}}{\color{oldcolor}}

% used here only to provide the XeLaTeX and BibTeX logos
\usepackage{hologo}

% if you use multirows in a table, include the multirow package
\usepackage{multirow}

% provides sidewaysfigure environment
\usepackage{rotating}

\usepackage{endnotes} % Inclui o pacote endnotes

% CUSTOM EPIGRAPH - BEGIN 
%%% https://tex.stackexchange.com/questions/193178/specific-epigraph-style
\usepackage{epigraph}
\renewcommand\textflush{flushright}
\makeatletter
\newlength\epitextskip
\pretocmd{\@epitext}{\em}{}{}
\apptocmd{\@epitext}{\em}{}{}
\patchcmd{\epigraph}{\@epitext{#1}\\}{\@epitext{#1}\\[\epitextskip]}{}{}
\makeatother
\setlength\epigraphrule{0pt}
\setlength\epitextskip{0.5ex}
\setlength\epigraphwidth{.7\textwidth}
% CUSTOM EPIGRAPH - END

% to use IPA symbols in unicode add
%\usepackage{fontspec}
%\newfontfamily\ipafont{CMU Serif}
%\newcommand{\ipa}[1]{{\ipafont #1}}
% and in the text you may use the \ipa{...} command passing the symbols in unicode

% LANGUAGE - BEGIN
% ARABIC
% for languages that use special fonts, you must provide the typeface that will be used
% \setotherlanguage{arabic}
% \newfontfamily\arabicfont[Script=Arabic]{Amiri}
% \newfontfamily\arabicfontsf[Script=Arabic]{Amiri}
% \newfontfamily\arabicfonttt[Script=Arabic]{Amiri}
%
% in the article, to add arabic text use: \textlang{arabic}{ ... }
%
% RUSSIAN
% for russian text we also need to define fonts with support for Cyrillic script
% \usepackage{fontspec}
% \setotherlanguage{russian}
% \newfontfamily\cyrillicfont{Times New Roman}
% \newfontfamily\cyrillicfontsf{Times New Roman}[Script=Cyrillic]
% \newfontfamily\cyrillicfonttt{Times New Roman}[Script=Cyrillic]
%
% in the text use \begin{russian} ... \end{russian}
% LANGUAGE - END

% EMOJIS - BEGIN
% to use emoticons in your manuscript
% https://stackoverflow.com/questions/190145/how-to-insert-emoticons-in-latex/57076064
% using font Symbola, which has full support
% the font may be downloaded at:
% https://dn-works.com/ufas/
% add to preamble:
% \newfontfamily\Symbola{Symbola}
% in the text use:
% {\Symbola }
% EMOJIS - END

% LABEL REFERENCE TO DESCRIPTIVE LIST - BEGIN
% reference itens in a descriptive list using their labels instead of numbers
% insert the code below in the preambule:
%\makeatletter
%\let\orgdescriptionlabel\descriptionlabel
%\renewcommand*{\descriptionlabel}[1]{%
%  \let\orglabel\label
%  \let\label\@gobble
%  \phantomsection
%  \edef\@currentlabel{#1\unskip}%
%  \let\label\orglabel
%  \orgdescriptionlabel{#1}%
%}
%\makeatother
%
% in your document, use as illustraded here:
%\begin{description}
%  \item[first\label{itm1}] this is only an example;
%  % ...  add more items
%\end{description}
% LABEL REFERENCE TO DESCRIPTIVE LIST - END


% add line numbers for submission
%\usepackage{lineno}
%\linenumbers


\begin{document}
\maketitle

\begin{figure}[htbp]
\centering
\begin{minipage}{.5\textwidth}
\includegraphics[width=\textwidth]{Fig1.png}
 \caption*{\fullcite{haidt2024geracao}.}
 \label{fig01}
\end{minipage}
\end{figure}

\begin{polyabstract}
\begin{abstract}
A resenha analisa a obra \citetitle{haidt2024geracao}, de Jonathan Haidt, sob a ótica dos impactos da tecnologia na formação humana e na linguagem. O objetivo é examinar a tese do autor sobre a drástica transição de uma infância baseada no brincar para uma infância baseada no celular, correlacionando-a à epidemia de transtornos mentais juvenis. O texto detalha os quatro danos fundamentais: privação social, privação de sono, fragmentação da atenção e vício. Por fim, conclui-se discutindo as tensões entre as propostas de banimento de smartphones e as diretrizes de letramento digital da BNCC, problematizando a aplicação de normas globais em cenários de desigualdade tecnológica e infraestrutural típicos do Sul Global.

\keywords{Cultura digital \sep Infância \sep Letramento digital \sep Saúde mental}
\end{abstract}

\pagebreak
\begin{english}
\begin{abstract}
The review analyzes Jonathan Haidt's \citetitle{haidt2024anxious} from the perspective of the impacts of technology on human formation and language. The objective is to examine the author's thesis on the drastic transition from a play-based childhood to a cell phone-based childhood, correlating it with the epidemic of youth mental disorders. The text details the four fundamental harms: social deprivation, sleep deprivation, attention fragmentation, and addiction. Finally, it concludes by discussing the tensions between proposals for banning smartphones and the BNCC's digital literacy guidelines, problematizing the application of global norms in scenarios of technological and infrastructural inequality typical of the Global South.

\keywords{Digital culture \sep Childhood \sep Smartphones \sep Digital literacy \sep Mental health}
\end{abstract}
\end{english}
% if there is another abstract, insert it here using the same scheme
\end{polyabstract}

\section{Introdução}\label{sec-intro}
\textit{A geração ansiosa: como a infância hiperconectada está causando uma epidemia de transtornos mentais}, de autoria de Jonathan Haidt, foi publicado originalmente em 2024. No Brasil, foi traduzido por Lígia Azevedo e editado pela Companhia das Letras. O autor é professor na Stern School of Business na Universidade de Nova York. Em 1992, obteve o título de doutor em psicologia social, área à qual vem dedicando suas pesquisas, com foco em psicologia moral e política.

\vspace{-2pt}
\section{Desenvolvimento e análise crítica}\label{sec-normas}
Já na Introdução, Haidt destaca a complexidade do tema em questão. Tendo suas conclusões sido alvo de controvérsias, afirma que pode ter errado em alguns pontos \cite[p. 25]{haidt2024geracao}. Além da Introdução e da Conclusão, o livro é dividido em quatro partes: I. A onda que varreu; II. O pano de fundo: o declínio da infância baseada no brincar; III. A grande reconfiguração: A ascensão da infância baseada no celular; IV. Ações coletivas para uma infância mais saudável. As quatro partes apresentadas são subdivididas em um total de doze capítulos que compõem a obra em questão. Nesta resenha, opto por apresentar os principais pontos de cada parte e, a seguir, tecer alguns comentários críticos às análises de Haidt.

No início da obra, Haidt convida o leitor a imaginar a reação dos pais caso seus filhos fossem recrutados para uma missão de ocupação humana no planeta Marte. Ao ponderar as implicações desse experimento, o autor reflete que a maioria dos responsáveis dificilmente daria a sua autorização. No entanto, embora a sociedade jamais permitisse a ida a Marte nesse contexto, agimos de forma análoga e, sem a devida cautela, permitimos que as crianças acessem um novo planeta (o ambiente digital) ao lhes entregar um smartphone. É neste mundo reconfigurado que elas crescem, um ambiente que Haidt lista como possuindo quatro características fundamentais: descorporificado, assíncrono, um-para-muitos e sem custódia de entrada e saída \cite[p. 18-19]{haidt2024geracao}.

Na Parte I, Haidt esclarece quem compõe a geração ansiosa que intitula a obra: trata-se da Geração Z (nascidos a partir de meados da década de 1990), a primeira a atravessar a puberdade com um smartphone nas mãos. O autor descreve o que chama de "Grande reconfiguração", período compreendido entre 2010 e 2015, quando os smartphones e as mídias sociais foram incorporados de forma mais proeminente na vida social, atingindo o desenvolvimento cognitivo e social desses jovens. O que ele denomina "onda gigante" é a epidemia de ansiedade, depressão e automutilação que acometeu as infâncias e as adolescências a partir desse marco temporal. O autor levanta a diversidade de estudos sobre os transtornos psíquicos para apontar como o crescimento foi exponencial neste período, o que o leva a sugerir que tal simultaneidade de eventos não é uma coincidência. Além disso, nota-se que essa onda gigante afetou as identidades de gênero de forma distinta, considerando a maneira particular como meninos e meninas exerciam suas performatividades nas redes e ambientes digitais.

Já na Parte II, Haidt levanta a hipótese de uma mudança estrutural, delineando a transição de uma infância baseada no brincar para uma infância baseada no celular (aqui tratado como sinônimo de smartphone). O autor estabelece um contraste direto entre os dois modelos: enquanto o brincar livre oferece riscos saudáveis (como o cair, o correr e o se esconder) que oportunizam o desenvolvimento de habilidades essenciais, os videogames e jogos online, ao serem considerados em sua forma mais ampla e corroborando as quatro características das relações mediadas, não suprem essa necessidade. Haidt enfatiza a importância do mundo físico, afirmando que "somos criaturas corporificadas; as crianças precisam aprender a administrar seu corpo no mundo físico antes de passar horas no virtual" \cite[p. 95]{haidt2024geracao}.

Na Parte III, o autor aprofunda a mecânica do adoecimento juvenil ao apresentar quatro prejuízos fundamentais no desenvolvimento de uma infância e adolescência baseadas no celular: privação social, privação de sono, atenção fragmentada e vício. Cumpre esclarecer que esses quatro elementos não configuram os transtornos mentais mencionados no título da obra, mas sim os fatores de risco e os mecanismos propulsores que deterioram a base psicológica dos jovens. São essas perdas contínuas que motivam o surgimento de patologias psíquicas, a exemplo da depressão, do Transtorno do Déficit de Atenção com Hiperatividade (TDAH) e do Transtorno de Ansiedade Generalizada (TAG), consolidando, assim, a epidemia de adoecimento mental que dá nome ao livro.

A relevância desses pontos foi constatada por uma observação que motivou a busca por esta obra: durante um mês de acompanhamento de turmas de aprendizagem (estudantes do Ensino Médio), chamou a atenção a baixa interação social entre os jovens que utilizavam fones de ouvido. Além disso, foi nítida a manifestação de atenção fragmentada, com alta dificuldade de entender comandos básicos e uma necessidade constante de checagem de notificações no celular.

Embora esses prejuízos delineados na resenha possam afetar também jovens, adultos e idosos, Haidt justifica seu foco nas infâncias e adolescências. Esse recorte permite ao autor fundamentar seus argumentos em dados e pesquisas específicas. Cumpre considerar que a criança e o adolescente se encontram em um estado de neuroplasticidade, o que confere aos impactos dos celulares nessas fases uma margem de desenvolvimento específica e mais crítica.

\vspace{-2pt}
\section{Conclusão}\label{sec-conduta}
A transposição das análises de Haidt para o contexto brasileiro exige mediações críticas. Um primeiro ponto é a universalidade da infância pressuposta pelo autor. O modelo de "superproteção" que ele critica é, muitas vezes, um privilégio de classe e raça. Em países do Sul Global, o celular não é apenas um ladrão de infância, mas frequentemente a única ferramenta de inclusão digital e lazer em territórios marcados pela violência urbana.

Além da questão sociológica, na minha perspectiva as propostas de Haidt encontram barreiras complexas na realidade educacional brasileira. Ao sugerir o banimento total de dispositivos nas escolas, percebe-se que o argumento do autor, quando transposto para o nosso contexto, acaba tensionando as diretrizes da própria Base Nacional Comum Curricular (BNCC), cuja competência 5 preconiza o domínio da cultura digital. Cria-se, assim, um paradoxo para o educador: como promover o letramento digital exigido pelo currículo excluindo o principal dispositivo de acesso disponível para estudantes de classes populares? A obra, portanto, é um diagnóstico necessário, mas sua aplicação exige uma tropicalização das estratégias, sob risco de confundir proteção com exclusão digital.


\vspace{-2pt}
\printbibliography\label{sec-bib}
% if the text is not in Portuguese, it might be necessary to use the code below instead to print the correct ABNT abbreviations [s.n.], [s.l.]
%\begin{portuguese}
%\printbibliography[title={Bibliography}]
%\end{portuguese} 

\vspace{-2pt}
\vspace{-1em}
\section*{\fontsize{9}{12}\selectfont Nota de transparência metodológica}
A elaboração desta resenha contou com apoio pontual de ferramentas de inteligência artificial generativa, utilizadas na revisão linguística e organização estrutural. A análise crítica e os juízos avaliativos são de inteira responsabilidade do autor, em conformidade com os princípios de ética científica ao qual se exige no âmbito da academia.

\vspace{-2pt}
\begin{dataavailability}
\txtdataavailability{nodata} % options: dataavailable, dataonly, databody, datanotav, nodata
\end{dataavailability}

\end{document}

